\subsection{Accuracy and Robustness Trade-offs}

The three depth-source modalities occupy distinct positions in the
accuracy--cost--robustness trade-off space:

\begin{itemize}
    \item \textbf{LiDAR + RGB + IMU}: Highest tracking accuracy and
    geometric fidelity, especially outdoors and under aggressive motion.
    LiDAR provides direct metric geometry that anchors the Gaussian map.
    However, LiDAR hardware is expensive ($$\,$$1k--$$\,$$10k+), bulky, and
    power-hungry.  Systems like Gaussian-LIC2 and RMGS-SLAM achieve
    sub-decimeter ATE in outdoor scenes where visual-only methods fail.

    \item \textbf{iToF + RGB + IMU}: Moderate accuracy, limited by iToF's
    short range ($<$10\,m) and ambient light sensitivity.  The extreme
    sparsity ($\leq$64 points) requires sophisticated depth completion.
    ToF-Splatting demonstrates that multi-frame integration can bridge the
    gap to RGB-D quality, but the approach remains indoor-focused.
    Hardware cost is minimal ($$\,$$5--$$\,$$50), targeting smartphones
    and AR/VR headsets.

    \item \textbf{Monocular Depth + RGB + IMU}: Accuracy depends heavily on
    the depth prior quality and IMU fusion strategy.  VIGS-SLAM achieves
    competitive results by tightly coupling IMU with learned dense
    correspondences.  HI-SLAM2 surpasses RGB-D methods on indoor benchmarks
    through monocular geometry priors + scale alignment.  However, outdoor
    large-scale scenes remain challenging due to depth ambiguity at range.
    Hardware cost is essentially zero (camera + IMU already in every
    smartphone).
\end{itemize}

\subsection{Depth Completion and Densification Strategies}

A common thread across all three categories is the need to \emph{densify}
the raw depth source before initializing Gaussians.  \Cref{tab:depth_comp}
summarizes the strategies:

\begin{table}[t]
\centering
\caption{Depth completion/densification strategies across modalities.}
\label{tab:depth_comp}
\scriptsize
\begin{tabular}{lll}
\toprule
\textbf{Modality} & \textbf{Strategy} & \textbf{Representative} \\
\midrule
LiDAR & Zero-shot depth completion network & Gaussian-LIC2 \\
LiDAR & Voxel-GPR interpolation & GS-LIVM \\
LiDAR & Voxel-PCA geometric priors & RMGS-SLAM \\
LiDAR & Depth propagation from dense maps & LV-GS SLAM \\
\midrule
iToF & Multi-frame integration (DoD) & ToF-Splatting \\
iToF & Monocular depth prior fusion & Fang et al. \\
\midrule
Mono & Foundation model depth (Depth Pro) & CGS-SLAM \\
Mono & OmniData depth + scale grid & HI-SLAM2 \\
Mono & DROID-SLAM dense disparity & VIGS-SLAM \\
Mono & Optical flow as geometry cue & GaussianFlow \\
Mono & VGGT feed-forward pointmaps & VGGT-GS SLAM \\
Mono & SfM + LFM fusion & VIM-GS \\
\bottomrule
\end{tabular}
\end{table}

\subsection{Real-Time Performance}

Real-time operation is a critical differentiator.  The key bottleneck is
Gaussian optimization, which involves differentiable rasterization and
gradient backpropagation at each keyframe.  Strategies to achieve real-time
performance include:

\begin{itemize}
    \item \textbf{Sliding-window optimization}: GS-LIVO optimizes only
    Gaussians within a temporal sliding window, reducing GPU memory and
    computation.
    \item \textbf{Custom CUDA kernels}: Gaussian-LIC, Gaussian-LIC2, LVI-GS,
    and VIGS-SLAM implement custom CUDA kernels for rasterization and
    optimization, avoiding Python overhead.
    \item \textbf{Submap management}: RMGS-SLAM and LV-GS SLAM use
    keyframe-based submaps with dynamic memory allocation.
    \item \textbf{Parallel tracking and mapping}: Most systems separate
    odometry and Gaussian mapping into parallel threads.
    \item \textbf{Embedded deployment}: GS-LIVO (Jetson Orin NX) and
    GS-LIVM (custom CUDA) demonstrate feasibility on resource-constrained
    platforms.
\end{itemize}

\subsection{Open-Source Code Availability}

Open-source code is critical for reproducibility and community adoption.
\Cref{tab:code} lists the systems with confirmed code releases:

\begin{table}[t]
\centering
\caption{Open-source code availability.}
\label{tab:code}
\scriptsize
\begin{tabular}{lcc}
\toprule
\textbf{System} & \textbf{GitHub URL} & \textbf{Category} \\
\midrule
Gaussian-LIC & APRIL-ZJU/Gaussian-LIC & LiDAR \\
GS-LIVM & xieyusen/GS-LIVM & LiDAR \\
Gaussian-LIC2 & xingxingzuo/gaussian\_lic2 & LiDAR \\
MonoGS & muskie82/MonoGS & Monocular \\
HI-SLAM2 & hi-slam2.github.io & Monocular \\
GaussianFlow & url-kaist/gaussianflow-slam & Monocular \\
ToF-Splatting & andreaconti.github.io & iToF \\
\bottomrule
\end{tabular}
\end{table}

\subsection{Loop Closure and Global Consistency}

Loop closure is essential for large-scale SLAM but is often omitted in
Gaussian SLAM systems to preserve runtime.  Only a few systems implement
explicit loop closure:

\begin{itemize}
    \item \textbf{RMGS-SLAM}: Gaussian-based GICP registration on the
    optimized map, followed by pose-graph optimization.
    \item \textbf{VIGS-SLAM}: Optical-flow-based loop detection with
    consistent Gaussian map updates after pose correction.
    \item \textbf{HI-SLAM2}: Sim(3) pose-graph BA with instant Gaussian
    deformation based on anchored keyframe updates.
    \item \textbf{VGGT-GS SLAM}: Gaussian-native alignment for
    loop-candidate verification before global Sim(3) pose-graph
    optimization.
\end{itemize}

The challenge is that after pose correction, the Gaussian map must be
updated to reflect the new keyframe poses---either by re-optimization
(expensive) or by explicit Gaussian deformation (efficient, as in HI-SLAM2).
